\cotefile{157-2}
\scanpage{18}
des foncteurs \inserted{$\unsure{c'|\varepsilon'}\colon\unsure{\mathscr{L}|\mathcal{L}}_{\unsure{m'n'|m'+n'|m',n'}}\unsure{\mathord{-}|\to}X^{\unsure{\checkmark|v}}$} dont le composé avec $\alpha'$ est le foncteur
\[
\beta\colon\varepsilon\longrightarrow X\qquad
\begin{aligned}a_0&\unsure{\mathrel{\mapstochar\relbar}|\longmapsto} x\\ b_0&\longmapsto z\end{aligned}
\tag*{$(*)$}
\]
et de $=$ la catégorie-source $\unsure{A/|A}$ \deleted{\unsure{est formée de la}} \inserted{s'interprète comme la}
sous-catégorie pleine de $\underline{\mathrm{Hom}}(\mathscr{L}_{\unsure{mn|m+n}},X)$ formée
des $c\colon\mathscr{L}_{mn}\mathord{-}X$ dont la restriction à
$\{a,b,c\}$ est le foncteur
\[
a\mathrel{\mapstochar\relbar}x,\ b\mathrel{\mapstochar\relbar}y,\ c\mathrel{\mapstochar\relbar}z\quad .
\]
Elle est contenue dans la sous-catégorie pleine
$A'$ des foncteurs dont la restriction à $\{a,b\}$ soit
$a\mathrel{\mapstochar\relbar}x,\ b\mapsto z$, \unsure{i.e.} dont le composé avec $\alpha$ soit
le foncteur $\beta\colon\varepsilon\mathord{-}X$ déjà envisagé ci-dessus.
Ainsi les foncteurs $f,g\colon A\rightrightarrows B$ se factorisent
par $A'$
\diagram{\begin{tikzcd}[ampersand replacement=\&,column sep=large]
 A \arrow[r,hook,no head] \arrow[drr,shift left=1,"f"] \arrow[drr,shift right=1,"g"'] \& A' \arrow[dr,shift left=1,"f'"] \arrow[dr,shift right=1,"g'"'] \& {}\\
 {} \& {} \& B
\end{tikzcd}},
foncteurs déduits de $\varphi,\psi$ ci-dessus, donnant
lieu \unsure{aux} triangles commutatifs \unsure{(3)}, et \deleted{\unsure{au}} foncteur
\deleted{\unsure{fonct.}} $\beta\colon\varepsilon\mathord{-}X$, de la façon dite. Pour
prouver que $f$ et $g$ sont homotopes, il suffit
de le prouver pour $f',g'$. Et pour ceci, il
\deleted{\unsure{revient au même}} \inserted{suffit} de montrer que $\varphi$ et $\psi$
sont homotopes en \unsure{tant} qu'objets de la
catégorie $\underline{\mathrm{Hom}}(\mathscr{L}_{m+n},\unsure{\mathscr{L}}_{\unsure{m'n'|m'+n'|m',n'}}\unsure{,};\alpha)$ formée
des foncteurs dont la restriction à $\varepsilon$ est $\alpha$.
% archive sheet number: 17
\scanpage{19}\mspage{17}
D'ailleurs, les deux assertions \inserted{(sur $f',g'$, sans préciser $X,\unsure{x|a},y,z$, et sur $\varphi,\psi,\alpha$)} sont en fait équivalentes,
comme on voit en faisant $X=\unsure{\mathscr{L}|\mathcal{L}}_{m+n},\ x=a,y=c,z=b$.
Ceci \unsure{se trouve|est contenu} dans le lemme \deleted{\unsure{préc.}} plus général

\uline{Lemme}. Soient $k$ et $l$ deux entiers. Alors la
sous-catégorie pleine de $\underline{\mathrm{Hom}}(\mathscr{L}_k,\mathscr{L}_l)$ formée des
foncteurs \deleted{\unsure{qui}} $\mathscr{L}_k\to\mathscr{L}_l$ appliquant source en source,
but en but est connexe. (NB Elle est non vide
\unsure{ssi} $k\geq l$).
% slanted
\margin{Dans les considérations qui suivent, \unsure{on considère|on utilise|où en vérité} en fait \unsure{tacitement que} les foncteurs envisagés sont \uline{croissants} sur les sommets (pour l'ordre de leur numérotation). Donc il faudrait écrire $\underline{\mathrm{Hom}}_{\mathrm{cr}}$ au lieu de $\underline{\mathrm{Hom}}$, dans les six énoncés qui suivent, y compris dans le lemme de connexité $\unsure{h'|k'}$ \deleted{\unsure{pour}} ci-contre.}

\deleted{\unsure{Un} tel foncteur, est décrit par une suite de sous-intervalles $I_\alpha$ de $\unsure{[1,2k+1]}$
l'intervalle d'entiers $[\unsure{1|2},2l]\simeq\mathrm{Fl}(\mathscr{L}_k)$, les
$I_\alpha$ de longueur $\lambda_\alpha$ paire, \inserted{mutuellement disjoints,} et avec $\sum\lambda_\alpha+l=k$,
et tel que le plus petit élément de $[1,2l]$ \unsure{compris dans} \inserted{qui ne soit dans aucun des}
$I_\alpha$ soit une flèche de même sens que la première flèche de $\mathscr{L}_l$.
Pour}
\[
\deleted{(I_\alpha)\in\unsure{\mathfrak{P}|\mathcal{P}}(\unsure{\mathfrak{P}|\mathcal{P}}(\mathscr{L}_k))\quad\unsure{;;}}
\]

On va décrire entièrement ces foncteurs \inserted{(\unsure{strt.|tacit.} croissants)}
$\mathscr{L}_k\xrightarrow{\varphi}\mathscr{L}_l$, en identifiant $\mathrm{Fl}(\mathscr{L}_k)$ \unsure{;|à} l'intervalle
d'entiers $[1,2k]$, et en associant à $\varphi$ la partie
\[
E_\varphi\subset\mathrm{Fl}(\mathscr{L}_k)\simeq[1,2k]
\]
\margin{$E_\varphi=\{u\in\mathrm{Fl}(\unsure{\mathscr{L}}_k)\mid\varphi(u)\ \text{une identité}\}$}
formée des $u\in\mathrm{Fl}(\mathscr{L}_k)$ \deleted{telles que} qui sont ``contractées''
par $\varphi$ en une identité. Si $E$ est un ensemble
d'entiers, on y introduit la notion de composantes
connexes (par la relation \inserted{d'équiv. de $E$} engendrée par la relation
de succession $y=x+1$). On voit aussitôt que
les composantes connexes de $E_\varphi$ \inserted{qui ne contiennent ni la première ni la dernière flèche de $\mathscr{L}_k$} sont paires, et
que
\[
\underbrace{\operatorname{card}\unsure{\mathrm{Fl}|\mathrm{El}}(\mathscr{L}_k)}_{2k}
=\operatorname{card}\lvert E_\varphi\rvert+\underbrace{\operatorname{card}\mathrm{Fl}(\mathscr{L}_{\unsure{l|a}})}_{2l}
\]
% archive sheet number: 18
\scanpage{20}\mspage{18}
\deleted{\unsure{donc}} \unsure{i.e.}
\[
\operatorname{card}E_\varphi=2(\deleted{\unsure{l}}k-\deleted{\unsure{k}}l)
\]
Il est immédiat qu'un foncteur $\varphi$ est surjectif
sur les objets et sur les flèches, donc identifie
$\unsure{\mathscr{L}|\mathcal{L}}_l$ à une catégorie quotient de $\mathscr{L}_k$ -- \deleted{\unsure{savoir}}
et la relation d'équivalence dans $\mathscr{L}_k$ définie
par $\varphi$ est celle obtenue en \deleted{\unsure{écrivant que les}} \inserted{identifiant}
tous les objets qui sont sources ou buts de
flèches d'une $=$ \unsure{composante|compos.te} connexe \inserted{$I$} de $E_\varphi$
\diagram{\begin{tikzcd}[ampersand replacement=\&,column sep=small]
\cdots \& {} \arrow[r,phantom] \& {} \arrow[l] \arrow[r] \& \bullet \& \bullet \arrow[l] \arrow[r] \& \bullet \& {} \arrow[l] \arrow[r] \& {} \& \cdots
\arrow[from=1-3,to=1-7,phantom,"{\underbrace{\hspace{6em}}_{I}}"']
\end{tikzcd}}
et les flèches dans cette composante à l'identité correspondante. Les conditions énoncées sur $E$
(indépendamment d'une $\varphi$ qui lui \unsure{donnerait}
\unsure{nai ssance|naissance}) assurent que la catégorie quotient
est isomorphe \inserted{(avec \uline{préservation} de source et de but)}
soit à $\mathscr{L}_l$, soit à la duale $\unsure{\breve{\mathscr{L}}|\check{\mathscr{L}}}_l=\mathscr{L}_l^{\unsure{\vee|\smile}}$,
suivant la direction de la première flèche de
\deleted{c'est le premier si $k$ et $l$ sont de même parité, et
le deuxième \deleted{\unsure{dans le}} cas contraire. (Car en
tous les cas, la catégorie quotient ``commence'' par
une flèche de $=$ direction que $\mathscr{L}_k$, laquelle
en dépend \unsure{pour} de la parité)}
$\overline{\mathscr{L}}_k$ (quotient de $\mathscr{L}_k$).
Cette ``direction \inserted{initiale}'' est la $=$ que celle pour $\mathscr{L}_k$ si
la première flèche \inserted{de $\mathscr{L}_k$} n'appartient pas à $E$, ou bien
si \deleted{sa comp.} elle y appartient et sa comp. conn.
dans $E$ est de cardinal pair. Dans le cas
inverse ($u_{\unsure{1|*}}\in E$ et sa comp. conn. dans $E$ est de
cardinal impair), \deleted{elle est de cont.} \inserted{la direction initiale de $\overline{\mathscr{L}}_k$ est opposée à celle de $\mathscr{L}_k$.}
% archive sheet number: 19
\scanpage{21}\mspage{19}
On voit donc que $E$ doit satisfaire encore à une
condition supplémentaire, pour que \inserted{le} $\overline{\unsure{\mathscr{L}|\mathcal{L}}}_k$ correspondant
soit \unsure{isom.|iso.} à $\mathscr{L}_l$ i.e. pour que $E$ provienne d'un
foncteur $\varphi$. Je résume toutes les conditions

\[
\operatorname{Card}E=2(k-l)\tag*{a)}
\]
b) Les comp. conn. de $E$ qui ne contiennent pas
de flèche initiale ni finale sont de card. pair.

c) Si $\mathscr{L}_k$ et $\mathscr{L}_l$ ont $=$ direction initiale
($k$ et $l$ de $=$ parité), alors l'éventuelle
composante de $E$ contenant la flèche
initiale doit être elle aussi de cardinal pair
(\unsure{absence signifiant card de 0|absence signifiant card de .}), équivaut à la $=$
condition pour l'éventuelle comp. conn. contenant la flèche finale de $\mathscr{L}_k$). Si $\mathscr{L}_k$ et $\mathscr{L}_l$
ont des directions initiales différentes, il
faut que la flèche initiale de $\mathscr{L}_k$ appartienne
à une comp. conn. de cardinal \uline{impair} de
$E$ (ce qui équivaut à idem pour la flèche
finale).

D'autre part, le foncteur $\varphi$ est uniquement déterminé par
la connaissance de $E$, car \deleted{ils sont} il y a \unsure{une|un}
\uline{unique} iso. $\overline{\mathscr{L}}_k\xrightarrow{\sim}\mathscr{L}_l$ compatible aux sources et
buts, à cause du

\uline{Lemme} (immédiat) : \deleted{\unsure{tout|tel}} autom. de $\mathscr{L}_l$ qui
est l'identité sur l'\deleted{\unsure{point}} \inserted{objet} source (ou sur l'objet
but) est l'identité partout.

Les foncteurs $\varphi,\psi$ de tantôt correspondent
au choix, l'un de $E_\varphi=\{1,2,k-1,k\}$, l'autre de
$E_\psi=\{1,2,i,i+1\}$ \deleted{\unsure{tous les deux}} (avec $2$ composantes connexes
% archive sheet number: 20
\scanpage{22}\mspage{20}
dans un cas, trois dans l'autre).

Dans le cas général, il y a lieu d'expliciter
la relation d'ordre de \unsure{la|l'} \deleted{\unsure{catégorie}} $\underline{\mathrm{Hom}}^{!}(\unsure{\mathscr{L}|\mathcal{L}}_k,\mathscr{L}_l)$
($!$ signifie : compatible aux source et but), en termes
des \deleted{$\unsure{E_\varphi}$} ensembles représentatifs $E_\varphi,E_\psi$.
Je ne vais pas expliciter cette relation d'ordre, mais
seulement une propriété de celle-ci qui nous sera
suffisante. \deleted{\unsure{Notons de}} Soit $\varphi_E$ le foncteur défini
par une partie $E\subset[1,2k]$. Tout d'abord :

\uline{Lemme \unsure{1|2|3}} (immédiat) : \inserted{Soient $E,E'\subset[1,2k]$ admissibles, et} soit $j\in[1,2k]$ tel que
$E\cap[1,j]=E'\cap[1,j]$. Alors $\varphi_E$ et $\varphi_{E'}$ coïncident sur \unsure{les}
le \deleted{sommet} \inserted{but} de $u_j$ ($j$.\unsure{ième} flèche de $\mathscr{L}_k$) et sur
tous les objets de $\mathscr{L}_k$ \deleted{\unsure{de rang}} antérieurs à celle-ci
(au sens de la numérotation des objets, et non de la
relation d'ordre de la catégorie $\unsure{\mathscr{L}}_k$ bien sûr). Dualement,
si $E\cap[j,2k]=E'\cap[j,2k]$, alors $\varphi_E$ et $\varphi_{E'}$ coïncident
sur la source de $u_j$, et sur tous les objets de $\mathscr{L}_k$ postérieurs
à celui-ci.

Considérons maintenant une \deleted{composante connexe} \inserted{partie connexe}
$I$ de $E$, \deleted{et soit $I'$} (pas forcément une composante
connexe), \unsure{soit} $I=[r,s]\subset[1,2k]$ ($r$ le plus petit élément de $I$,
$s$ son \deleted{élément} plus grand élément.) et supposons que

a) $\operatorname{card}I$ est pair (i.e. $s-r$ est pair)

b) $r\ne1$, et $r-1\notin E$ (i.e. $r-1\in[1,2k]\setminus E$)

et soit \deleted{$\unsure{E'=}$}
\[
E'=(E\setminus\{s\})\cup\{r-1\}=(E\setminus I)\cup\overbrace{[r-1,s-1]}^{I'}
\]
\unsure{dérivant|déduit} de $I$ en ``translatant'' $I$ sur la gauche d'un cran.
% archive sheet number: 21
\scanpage{23}\mspage{21}
Il est immédiat que si $E$ est admissible, $E'$ aussi (et
inversement). Ceci dit, on a le

\uline{Lemme 2} \quad $E$ et $E'$ étant comme ci-dessus, il existe un
\unsure{hom|Hom} \deleted{\unsure{fonct.}} de foncteurs entre $\varphi_E$ et $\varphi_{E'}$ (je ne dis
pas dans quel sens va la flèche !) \deleted{\unsure{$\leftarrow$}}.

\uline{Cor} \quad Énoncé dual, quand $E'$ est déduit de $E$ en translatant
à droite \inserted{d'un cran} un intervalle \inserted{connexe} de longueur paire \inserted{$I=[r,s]$} de $E$, \deleted{tel que l'on ait}
$s\ne 2k$ et $s+1\notin E$.

En effet, dans ce cas $E$ est déduit de $E'$ par une \unsure{op.} du
type explicité dans le lemme 2.

La connexité de la catégorie $\underline{\mathrm{Hom}}^{!}(\unsure{\mathscr{L}|\mathcal{L}}_k,\mathscr{L}_l)$ est
maintenant claire. Si $k$ et $l$ ont même parité,
en poussant les comp. connexes de $E$ vers la gauche,
on met $E$ sous la forme ``normale''
\[
E_0=[1,2(k-l)]\quad .
\]
Si $k$ et $l$ sont de parité inverse, on \deleted{\unsure{ne}} ne peut
pousser \deleted{tout} la comp. conn. de $E$ qui contient
l'élément dernier $2k$, car elle est de cardinal impair.
Il faut laisser $2k$ dans $E$ et pousser le reste, et
on ramène $E$ à la forme normale
\[
E_0=[1,2(k-l)-1]\cup\{2k\}\quad ,
\]
cqfd !

\deleted{\unsure{§}} Ce qui précède justifie notre raisonnement,
qui donnait par composition des chemins des
applications de composition bien définies
\[
\underbrace{\operatorname{Hom}_{\unsure{\Sigma|\overline{Z}|\overline{\Sigma}|Z}}(x,y)}_{\substack{\Vert\\\pi_0\underline{\mathrm{Ch}}_\Sigma(x,y)}}
\times\underbrace{\operatorname{Hom}_{\Sigma}(y,z)}_{\substack{\Vert\\\pi_0\underline{\mathrm{Ch}}_\Sigma(y,z)}}
\unsure{\to|\mathord{-}}
\underbrace{\operatorname{Hom}_{\Sigma}(x,z)}_{\substack{\Vert\\\pi_0\underline{\mathrm{Ch}}_\Sigma(x,z)}}
\]
% archive sheet number: 22
\scanpage{24}\mspage{22}
Je dis que cette composition est associative. Cela
résulte du fait que si on \unsure{a} trois chemins (d'indices
\uline{pairs} !)
\diagram{\begin{tikzcd}[ampersand replacement=\&]
x \arrow[r,no head,"c"] \& y \arrow[r,no head,"c'"] \& \deleted{\unsure{t}}z \arrow[r,"c''"] \& \deleted{\unsure{z}}t
\end{tikzcd}},
on a \unsure{$^{*}$} (sans avoir à effectuer des \unsure{homotopies}\ldots)
associativité stricte
\[
c''\mathbin{.}(c'\mathbin{.}c)=(c''\mathbin{.}c')\mathbin{.}c\quad .
\]
Je dis qu'il y a des identités. On définit
\[
1_{\unsure{x|z}}\in\operatorname{Hom}_{\unsure{\Sigma|\overline{Z}|\overline{\Sigma}|Z}}(x,x)
\]
par le chemin \deleted{de longueur} \inserted{d'indice} $0$ (qui est pair !)
i.e. par le diagramme réduit à $x$, sans flèche
(ou si on veut, en interprétant le diagramme comme un
foncteur $e\to X$, il y a la flèche identique
sous-entendue). Nos raisonnements \unsure{n'excluaient|n'excluent}
pas le cas $n=0$ ou $m=0$, et il est \deleted{évident} \inserted{tautologique} que
pour un chemin
\diagram{\begin{tikzcd}[ampersand replacement=\&]x\arrow[r,no head,"c"]\&y\end{tikzcd}}
on a
\[
c\mathbin{.}1_x=1_y\mathbin{.}c=c\quad .
\]
Donc on a bien défini une catégorie, que je
vais appeler $X^{\deleted{\unsure{*}}}_{\Sigma}$. Je vais définir un
foncteur canonique
\[
\deleted{\unsure{Q|\Theta}}\unsure{\varphi|\Theta|Q}\colon X\longrightarrow X_{\Sigma}
\]
qui sera l'identité sur les objets, il faut le
définir sur les flèches
\[
\operatorname{Hom}_X(x,y)\unsure{\mathord{-}|\to}\operatorname{Hom}_{X_{\Sigma}}(x,y)=\pi_0\mathrm{Ch}_{\Sigma}(x,y)
\]
Malheureusement, une flèche
\diagram{\begin{tikzcd}[ampersand replacement=\&]x\arrow[r,no head,"u"]\&y\end{tikzcd}}
% archive sheet number: 23
\scanpage{25}\mspage{-- 23}
ne peut s'interpréter directement comme un chemin
\unsure{au|en} notre sens, il faut payer nos fautes ! Mais
on va à une telle flèche associer le chemin
\diagram{\begin{tikzcd}[ampersand replacement=\&,column sep=1.15em]
x \arrow[r] \& y \arrow[r,<->,"{\mathrm{id}}"] \& y \arrow[r,"{\mathrm{id}}"] \& y \arrow[r,<->,"{\mathrm{id}}"] \& y
\end{tikzcd}},
\deleted{de} \inserted{élément de} $\underline{\unsure{\mathrm{Ch}|\mathscr{L}}}_{\unsure{\Sigma|Z},2}(x,y)$, ce qui définit l'application \unsure{cherchée|du début}.
Il faut vérifier la compatibilité avec les identités
et avec la composition. À la flèche $1_x$ \inserted{$\colon x\mathord{-}x$} est associé
l'élément de $\unsure{\mathscr{L}|\mathcal{L}}_2(x,x)$
\diagram{\begin{tikzcd}[ampersand replacement=\&,column sep=1.15em]
x \arrow[r,no head,"{\mathrm{id}}"] \& x \arrow[r,<-,"{\mathrm{id}}"] \& x \arrow[r,"{\mathrm{id}}"] \& x \arrow[r,<-,"{\mathrm{id}}"] \& x
\end{tikzcd}},
Il faut prouver que dans $\pi_0(\underline{\mathrm{Ch}}_{\Sigma,2}(x,x))$ il devient
égal à l'élément déduit de $1_x\in\mathscr{L}_{\Sigma,0}(x,x)$,
par le foncteur de transition. Mais cet élément
n'est autre justement que
\diagram{\begin{tikzcd}[ampersand replacement=\&,column sep=1.15em]
x \arrow[r,"{\mathrm{id}}"] \& x \arrow[r,<-,"{\mathrm{id}}"] \& x \arrow[r,"{\mathrm{id}}"] \& x \arrow[r,<-,"{\mathrm{id}}"] \& x
\end{tikzcd}},
c'est le $=$ ! (sans avoir $=$ à passer au $\pi_0$) !

Pour la compatibilité avec la composition
\diagram{\begin{tikzcd}[ampersand replacement=\&,column sep=1.15em]
x \arrow[r,"{u}"] \& y \arrow[r,no head,"{v}"] \& z \arrow[r,<-] \& {}
\end{tikzcd}}
On a d'une part
\diagram{\begin{tikzcd}[ampersand replacement=\&,column sep=1.15em]
x \arrow[r,no head,"{vu}"] \& z \arrow[r,<-,"{\mathrm{id}}"] \& z \arrow[r,no head,"{\mathrm{id}}"] \& z \arrow[r,<-,"{\mathrm{id}}"] \& z
\end{tikzcd}}
d'\deleted{$x\to y\mathord{-}$}autre part
\diagram{\begin{tikzcd}[ampersand replacement=\&,column sep=1.15em]
x \arrow[r,"{u}"] \& y \arrow[r,<-,"{\mathrm{id}}"] \& y \arrow[r,"{\mathrm{id}}"] \& y \arrow[r,<-,"{\mathrm{id}}"] \& y \arrow[r,"{v}"] \& z \arrow[r,<-,"{\mathrm{id}}"] \& z \arrow[r,"{\mathrm{id}}"] \& z \arrow[r,<-,"{\mathrm{id}}"] \& z
\end{tikzcd}}
on complète le premier chemin par dégénérescence
aux deux bouts, pour le comparer au deuxième,
en tant qu'éléments de $\underline{\unsure{\mathrm{Ch}|\mathscr{L}}}_4(x,z)$ ;
\diagram{\begin{tikzcd}[ampersand replacement=\&,column sep=1.15em]
x \arrow[r,"{\mathrm{id}}"] \& x \arrow[r,<-,"{\mathrm{id}}"] \& x \arrow[r,no head,"{vu}"] \& z \arrow[r,<->,"{\mathrm{id}}"] \& z \arrow[r,<->,"{\mathrm{id}}"] \& z \arrow[r,<-,"{\mathrm{id}}"] \& z \arrow[r,"{\mathrm{id}}"] \& z \arrow[r,<-,"{\mathrm{id}}"] \& z
\end{tikzcd}}
% archive sheet number: 24
\scanpage{26}\mspage{24}
il faut prouver que les deux sont homotopes. On
a donc deux foncteurs
\diagram{\begin{tikzcd}[ampersand replacement=\&,column sep=1.5em]
(*)\quad f,g\colon
\underbrace{\overbrace{\underline{\mathrm{Fl}}(X;x,y)}^{\text{cat. discr.}}\times
\overbrace{\underline{\mathrm{Fl}}(X,y,z)}^{\text{cat. discr.}}}_{A}
\arrow[r,no head,shift left=1,"f"]\arrow[r,no head,shift right=1,"g"']
\& \underbrace{\unsure{\mathrm{Ch}|\mathscr{L}}_{4,\unsure{\Sigma|Z}}(X;x,z)}_{B}
\end{tikzcd}},
on veut prouver qu'ils sont homotopes. Pour ceci,
on les regarde \deleted{comme tantôt} \inserted{(comme tantôt)} comme ``transposés'' de deux
foncteurs
\diagram{\begin{tikzcd}[ampersand replacement=\&,column sep=large]
\mathord{\tikz[baseline=0pt,x=1ex,y=1ex]{\draw[line width=.09ex,line join=round] (0,0) -- (.95,1.30) -- (1.52,0) -- cycle; \draw[line width=.09ex] (.32,0) -- (1.02,1.12);}}_1\unsure{\amalg|\sqcup}\mathord{\tikz[baseline=0pt,x=1ex,y=1ex]{\draw[line width=.09ex,line join=round] (0,0) -- (.95,1.30) -- (1.52,0) -- cycle; \draw[line width=.09ex] (.32,0) -- (1.02,1.12);}}_1\simeq\mathord{\tikz[baseline=0pt,x=1ex,y=1ex]{\draw[line width=.09ex,line join=round] (0,0) -- (.95,1.30) -- (1.52,0) -- cycle; \draw[line width=.09ex] (.32,0) -- (1.02,1.12);}}_2
\& \unsure{\mathscr{L}|\mathcal{L}}_4 \arrow[l,shift right=1,"\varphi"']\arrow[l,shift left=1,"\psi"]\\
{} \& \varepsilon\arrow[u,"\alpha'"']\arrow[ul,"\alpha"']
\end{tikzcd}},
vu que les catégories \inserted{$A,B$} des deux \deleted{\unsure{buts}} de $(*)$
s'interprètent comme des sous-catégories pleines
de $\underline{\mathrm{Hom}}(\mathord{\tikz[baseline=0pt,x=1ex,y=1ex]{\draw[line width=.09ex,line join=round] (0,0) -- (.95,1.30) -- (1.52,0) -- cycle; \draw[line width=.09ex] (.32,0) -- (1.02,1.12);}}_2,X)$ et $\underline{\mathrm{Hom}}(\mathscr{L}_4,X)$ respectivement ;
à savoir des sous-catégories pleines formées
des foncteurs dont le composé avec $\alpha$ resp.$^{\mathrm{t}}$
$\alpha'$ est donné par $x$ et par $z$, et de plus
\deleted{\unsure{du fait de}}\unsure{(pour} $A$) dont la valeur en l'objet médian de
$\mathord{\tikz[baseline=0pt,x=1ex,y=1ex]{\draw[line width=.09ex,line join=round] (0,0) -- (.95,1.30) -- (1.52,0) -- cycle; \draw[line width=.09ex] (.32,0) -- (1.02,1.12);}}_2$ soit $y$. On peut encore remplacer $A$
par la ss-catégorie pleine plus grande \inserted{$A'$} de
$\underline{\mathrm{Hom}}(\mathord{\tikz[baseline=0pt,x=1ex,y=1ex]{\draw[line width=.09ex,line join=round] (0,0) -- (.95,1.30) -- (1.52,0) -- cycle; \draw[line width=.09ex] (.32,0) -- (1.02,1.12);}}_2,X)$, obtenue en ignorant cette
dernière condition, et il \deleted{\unsure{est évident}} \inserted{suffit de}
prouver que les deux foncteurs ``plus gros''
\diagram{\begin{tikzcd}[ampersand replacement=\&]
A'\arrow[r,shift left=1,"f'"]\arrow[r,shift right=1,"{g^{\unsure{\prime|=}}}"']\& B
\end{tikzcd}}
sont homotopes. Et pour ceci, il suffit de
% archive sheet number: 25
\scanpage{27}\mspage{25}
voir encore que $\varphi$ et $\psi$ sont homotopes dans la
catégorie des foncteurs $\unsure{\mathscr{L}|\mathcal{L}}_4\unsure{\mathord{-}|\to}\mathord{\tikz[baseline=0pt,x=1ex,y=1ex]{\draw[line width=.09ex,line join=round] (0,0) -- (.95,1.30) -- (1.52,0) -- cycle; \draw[line width=.09ex] (.32,0) -- (1.02,1.12);}}_2$ \deleted{dont la} qui
sont compatibles avec la source, et le but.

On voit donc que notre lemme d'homotopie
de tantôt n'était pas assez général. Il faudrait
définir les \uline{catégories \unsure{intervalles|intervalle} généralisées}
de longueur $n$ quelc. ($n=\mathrm{nb}$ de flèches bout
à bout), une catégorie d'int. gén. \inserted{\uline{\unsure{standart|standard}}} de long.
$n$ a pour \deleted{\unsure{objets les}} \inserted{ens.} d'objets l'int. d'entiers
\[
[0,\deleted{\unsure{2}}n],
\]
et pour tout couple d'entiers consécutifs
$i-1,i\deleted{\unsure{+1}}$ \deleted{\unsure{qui sont}} \unsure{en} objets ($\unsure{0|1}\leq i\leq n$), il y a
exactement une flèche qui relie $i,i+1$. \deleted{\unsure{Et}}

On affecte l'intervalle $]i,i+1]$ du signe $+$ si
$i\deleted{\leftarrow}\to i+1$ (\inserted{flèche en} sens direct), du signe \unsure{$i$|$-$|$\div$} si
$i\leftarrow i+1$ (flèche en sens rétrograde) -- donc
les catégories \unsure{standart|standard} sont définies, pour $n$ fixé,
par un système de signes $+$ ou $-$, dont on affecte
les $n$ ``\unsure{tronçons}'' consécutifs $\deleted{\unsure{A[i,i+1]}}\inserted{A_i=}[i-1,i]$ ($1\leq i\leq n$)
de l'intervalle d'entiers $[\unsure{0|1},n]$. Il y en a
donc exactement $2^n$.

On s'attend à ce que l'énoncé suivant soit
vrai :

\uline{Lemme} \quad Soient $I,I'$ deux intervalles généralisés,
munis d'objets source et objet but (de sorte que
$I,I'$ sont \uline{canoniquement} isom. aux int. gén.
\unsure{standart|standard} correspondants). Soit $\underline{\mathrm{Hom}}^!(I,I')$
% archive sheet number: 26
\scanpage{28}\mspage{26}
la cat des foncteurs de $I$ dans $I'$, compatibles
avec source et but. Si cette catégorie est $\ne\varnothing$,
elle est $0$-connexe (et, je présume, même
contractile ?).

Je vais admettre provisoirement qu'il en est ainsi,
\inserted{du moins} dans le cas particulier de $\unsure{\mathscr{L}|\mathcal{L}}_4$ (intervalle gén. de long. $8$)
et $\mathord{\tikz[baseline=0pt,x=1ex,y=1ex]{\draw[line width=.09ex,line join=round] (0,0) -- (.95,1.30) -- (1.52,0) -- cycle; \draw[line width=.09ex] (.32,0) -- (1.02,1.12);}}_2$. Cela donne donc un foncteur
\[
\Theta\colon X\longrightarrow X_{\unsure{\Sigma|Z}}\quad .
\]
On a gagné, car la propriété universelle de ce
foncteur est pratiquement \unsure{évidente|immédiate}. \deleted{Cela établit} \inserted{Mais} j'ai
oublié de vérifier que $\Theta$ transforme \deleted{objets} \inserted{\unsure{élts|él.}} de $\Sigma$
\deleted{dans $X_\Sigma$}
en flèches inversibles. Soit donc
\diagram{\begin{tikzcd}[ampersand replacement=\&,column sep=1.15em]
x \arrow[r,"{u}"] \& y \arrow[r,<-] \& {}
\end{tikzcd}}
avec $u\in\Sigma$, on va définir un chemin en sens
inverse
\diagram{\begin{tikzcd}[ampersand replacement=\&,column sep=1.15em]
y \arrow[r,no head,"{c}"] \& x
\end{tikzcd}}
qui définira dans $X_\Sigma$ une flèche inverse de $u$.
L'idée naturelle est de prendre le ``chemin''
\diagram{\begin{tikzcd}[ampersand replacement=\&]\deleted{\unsure{x|y}}\&\deleted{\unsure{z|x}}\arrow[l,"{\deleted{\unsure{u}}}"]\end{tikzcd}}
\diagram{\begin{tikzcd}[ampersand replacement=\&,column sep=1.15em]
y \arrow[r,<-,"{u}"] \& x
\end{tikzcd}}
de $y$ dans $x$, mais ce n'est pas un chemin au sens
de nos définitions, on va le compléter en
\diagram{\begin{tikzcd}[ampersand replacement=\&,column sep=1.15em]
c\colon\quad y \arrow[r,"{\mathrm{id}}"] \& y \arrow[r,<-,"{u}"] \& x \arrow[r,no head,"{\mathrm{id}}"] \& x \arrow[r,<-,"{\mathrm{id}}"] \& x
\end{tikzcd}}
Je dis que $c$ fait l'affaire. on doit prouver
deux relations d'homotopie, de \deleted{$\unsure{c\ast u}$} $c.\bar{u}$ et
de $\bar{u}.c$ avec les chemins\deleted{\unsure{de}} identiques \unsure{en|au} $x$ de long.
$8$ (indice $4$ i.e. dans $\underline{\mathrm{Ch}}_{\Sigma,4}(X;x,\unsure{x|X})$ resp $\underline{\mathrm{Ch}}_{\Sigma,4}(X;y,y)$)
% archive sheet number: 27
\scanpage{29}\mspage{27}
où $\bar{u}$ est le chemin associé à $u$
\diagram{\begin{tikzcd}[ampersand replacement=\&,column sep=1.15em]
x \arrow[r,no head,"{u}"] \& y \arrow[r,<-,"{\mathrm{id}}"] \& y \arrow[r,no head,"{\mathrm{id}}"] \& y \arrow[r,<-,"{\mathrm{id}}"] \& y
\end{tikzcd}}.
On a donc à regarder
\diagram{\begin{tikzcd}[ampersand replacement=\&,column sep=1.15em]
x \arrow[r,no head,"{u}"] \& y \arrow[r,<-,"{\mathrm{id}}"] \& y \arrow[r,no head,"{\mathrm{id}}"] \& y \arrow[r,<-,"{\mathrm{id}}"] \& y \arrow[r,no head,"{\mathrm{id}}"] \& y \arrow[r,<-,"{u}"] \& x \arrow[r,no head,"{\mathrm{id}}"] \& x \arrow[r,<-,"{\mathrm{id}}"] \& x
\end{tikzcd}}
et voir qu'il est homotope au chemin trivial
\diagram{\begin{tikzcd}[ampersand replacement=\&,column sep=1em]
x\arrow[r,"\mathrm{id}"]\&x\arrow[r,<-,"\mathrm{id}"]\&{}\&\cdots\&x\arrow[r,no head,"\mathrm{id}"]\&\unsure{X|x}\arrow[r,<-,"\mathrm{id}"]\&x
\end{tikzcd}}
(des identités partout), et item pour
\diagram{\begin{tikzcd}[ampersand replacement=\&,column sep=1.15em]
y \arrow[r,no head,"{\mathrm{id}}"] \& y \arrow[r,<-,"{u}"] \& x \arrow[r,no head,"{\mathrm{id}}"] \& x \arrow[r,<-,"{\mathrm{id}}"] \& x \arrow[r,no head,"{u}"] \& y \arrow[r,<-,"{\mathrm{id}}"] \& y \arrow[r,"{\mathrm{id}}"] \& y \arrow[r,<-,"{\mathrm{id}}"] \& y
\end{tikzcd}}
et le chemin trivial en $y$. Cette fois, on
a \inserted{encore} deux foncteurs
\diagram{\begin{tikzcd}[ampersand replacement=\&,column sep=large]
\mathord{\tikz[baseline=0pt,x=1ex,y=1ex]{\draw[line width=.09ex,line join=round] (0,0) -- (.95,1.30) -- (1.52,0) -- cycle; \draw[line width=.09ex] (.32,0) -- (1.02,1.12);}}_1 \& \unsure{\mathscr{L}|\mathcal{L}}_4\arrow[l,shift right=1,"\varphi"']\arrow[l,shift left=1,"\psi"]\\
{} \& \varepsilon\arrow[u,"\alpha"']\arrow[ul,"{\unsure{\alpha'|\alpha_1}}"']
\end{tikzcd}}
\deleted{\unsure{transfor}} compatibles avec \inserted{objets} sources et \inserted{objets} buts, et on doit
prouver qu'ils sont homotopes dans $\underline{\mathrm{Hom}}^{!}(\mathscr{L}_4,\mathord{\tikz[baseline=0pt,x=1ex,y=1ex]{\draw[line width=.09ex,line join=round] (0,0) -- (.95,1.30) -- (1.52,0) -- cycle; \draw[line width=.09ex] (.32,0) -- (1.02,1.12);}}_1)$.
Cela doit donc résulter du même lemme
général. On a \unsure{gagné|``gagné''} -- par une interprétation
simple de $X\unsure{\Sigma|Z}^{-1}$, et en particulier des
$\operatorname{Hom}_{X\Sigma^{-1}}(x,y)$, \inserted{et de leur composition,} en termes de la géométrie
des espaces de chemins :
\[
\left\{\begin{aligned}
&\operatorname{Hom}_{X\Sigma^{-1}}(x,y)\simeq\pi_0\bigl(\underline{\mathrm{Ch}}_\Sigma(X;x,y)\bigr)\\
&\text{Composition des hom dans }X\Sigma^{-1}\text{ correspond à}\\
&\qquad\text{la composition des $\Sigma$-chemins dans $X$}
\end{aligned}\right.
\]
% archive sheet number: 28
\scanpage{30}\mspage{28}
3. \uline{Passage aux catégories de diagrammes}

Nous \deleted{allons} convenons de noter la catégorie $X$
$\unsure{M|\mathcal{M}}$, \unsure{car} par la suite elle jouera le rôle d'une
``catégorie de modèles''. On désignera par contre
par $X$ les ``catégories d'indices'' sur lesquelles
on veut construire les \deleted{\unsure{fib. de}} ``catégories de
diagrammes'' de type $X$
\[
M(X)=\underline{\mathrm{Hom}}(X,M)\quad .
\]
On se donne donc
\diagram{\begin{tikzcd}[ampersand replacement=\&,row sep=1ex,column sep=0.4em]
\unsure{\Sigma|Z}\arrow[d] \& {}\subset\mathrm{Fl}(M), \& \Sigma\supset\text{ens. des flèches identiques},\\
\text{\uline{quasi-isomorphismes}}\&{}\&{}
\end{tikzcd}}
et pour \unsure{tte} catégorie $X$, on définit
\[
\Sigma_{\inserted{X}}\deleted{(\unsure{X})}\subset\mathrm{Fl}\,M(X)
\]
\unsure{ainsi} : Si $F,G\in\mathrm{Ob}\,M(X)$ ($F,G\colon X\rightrightarrows M$),
une flèche $u\colon F\unsure{\mathord{-}|\to}G$ (\unsure{hom|Hom} de foncteurs)
est dans $\Sigma_X\deleted{(\unsure{X})}$ ssi $\forall\,x\in X$, $u(x)\colon F(x)\mathord{-}G(x)$
est dans $\Sigma$ (``critère \inserted{de quasi-isomorphie} fibre par fibre'', pour
\inserted{l'hom. entre} les copréfaisceaux $F,\unsure{G|K}$ envisagés).

Soient
\[
a,b\in\mathrm{Ob}\,M
\]
et considérons les foncteurs constants correspondants
sur $X$
\[
a_X,b_X\in\mathrm{Ob}\,M(X),
\]
% archive sheet number: 29
\scanpage{31}\mspage{29}
On se propose d'appliquer \unsure{la|le} \deleted{``calcul'' préc.} \inserted{\unsure{théorie|théorème} du} \S\,2
au ``calcul'' \unsure{des|de}
\[
\operatorname{Hom}_{\unsure{\Sigma|Z}_X}(a_X,b_X),
\]
on a donc une \unsure{bij.} can.
\[
\boxed{\operatorname{Hom}_{\Sigma_X}(a_X,b_X)\simeq\pi_0\underline{\mathrm{Ch}}_{\deleted{\unsure{M(X)}}\Sigma_X}\bigl((\unsure{M|\mathcal{M}}(X);a_X,b_X\bigr).}
\]
On se propose d'expliciter la catégorie des
chemins du second membre, On a, abstraction
faite des données de $\Sigma,a,b$ dans $M$, la
\deleted{$\unsure{\underline{\mathrm{Ch}}_{\Sigma_X}}$}
expression pour la catégorie de \uline{tous} les
chemins dans $M(X)$
\[
\begin{aligned}
\underline{\mathrm{Ch}}(M(X))
&\simeq\underline{\mathrm{Hom}}\bigl(
\underset{\substack{\text{pro-cat.}\\(\unsure{\mathscr{L}|\mathcal{L}}_n)_{n\in\unsure{\mathbb{N}|N}}\ \\ \deleted{\unsure{n\ \mathrm{pair}}}}}{\unsure{\mathbb{L}|\mathscr{L}}},
\underbrace{M(X)}_{\underline{\mathrm{Hom}}(X,M)}\bigr)\quad )\\
&\overset{\mathrm{def}}{\simeq}\varinjlim_n
\underrightarrow{\underline{\mathrm{Hom}}}\bigl(\mathscr{L}_n,\underbrace{\overbrace{M(X)}}\bigr)\\
&\simeq\varinjlim_n\underbrace{\underline{\mathrm{Hom}}(\mathscr{L}_n\times X,M)}_{
\simeq\underline{\mathrm{Hom}}\bigl(X,
\underbrace{\underline{\mathrm{Hom}}(\mathscr{L}_n,M)}_{\substack{\Vert\,\mathrm{def}\\\underline{\mathrm{Ch}}_n(M)}}\bigr)}\\
&\simeq\varinjlim_n\underline{\mathrm{Hom}}\bigl(X,\underline{\mathrm{Ch}}_n(M)\bigr)\\
&\simeq\underline{\mathrm{Hom}}^!\bigl(X,\underline{\mathrm{Ch}}(M)\bigr)
\end{aligned}
\]
où le signe $!$ signifie qu'on prend la sous-catégorie pleine des foncteurs \deleted{$\unsure{X}$}
\[
X\unsure{\mathord{-}|\to}\underline{\mathrm{Ch}}(M)
\]
% archive sheet number: 30
\scanpage{32}\mspage{30}
qui se factorisent par une des sous-catégories
pleines $\underline{\deleted{\unsure{\mathscr{L}|\mathcal{L}}}\mathrm{Ch}}_n(\unsure{M|\mathcal{M}})$. On peut dire aussi que
si on considère le système inductif des
cat. $\underline{\mathrm{Ch}}_n(M)$ comme un ind-objet
de Cat, \deleted{\unsure{disons}} qu'on va noter
\[
\underline{\underline{\mathrm{Ch}}}_{\deleted{\unsure{n|\infty}}}(M)
\qquad\text{ind. objet de Cat},
\]
on a
\[
\underline{\mathrm{Ch}}(M(X))\overset{\sim}{\simeq}
\underline{\mathrm{Hom}}\bigl(X,\underline{\underline{\mathrm{Ch}}}_{\deleted{\unsure{n|\infty}}}(M)\bigr)\quad .
\]
Si on en vient \deleted{\unsure{aux}} à $\unsure{\Sigma|Z},\Sigma_X$, On trouve
aussitôt qu'un chemin dans $M(X)$ est un
$\Sigma_X$-chemin, ssi \deleted{\unsure{tous ses objets}} \inserted{le foncteur} correspondant
$X\unsure{\mathord{-}|\to}\underline{\mathrm{Ch}}(M)$ se factorise par $\underline{\mathrm{Ch}}_\Sigma(M)$.

Ainsi on trouve
\[
\underline{\mathrm{Ch}}_{\Sigma_X}(M(X))\simeq\underline{\mathrm{Hom}}^!(X,\underline{\mathrm{Ch}}_\Sigma(M)),
\]
où le signe $!$ a le même sens que plus bas.
Si on introduit le ind-objet $\underline{\underline{\mathrm{Ch}}}_\Sigma(M)$
\unsure{défini|définie} par les $\underline{\mathrm{Ch}}_{\Sigma,n}(M)$, on trouve donc
aussi
\[
\underline{\mathrm{Ch}}_{\Sigma_X}(M(X))\xrightarrow{\approx}
\underline{\mathrm{Hom}}\bigl(X,\underline{\underline{\mathrm{Ch}}}_\Sigma(M)\bigr)\quad .
\]
% archive sheet number: 31
\scanpage{33}
\mspage{31}

Soit $c$ un \deleted{\illegible{1}} chemin dans $\unsure{M|\mathcal{M}}(X)$, et
\[
  \deleted{\unsure{c\colon}} X\unsure{\mathord{-}|\longrightarrow}\underline{\mathrm{Ch}}.(\deleted{\unsure{X}}M)
\]
le foncteur qu'il définit, \unsure{composons-le} avec les foncteurs source-but \inserted{$\underline{\mathrm{Ch}}(\deleted{\unsure{X}}M)\unsure{\mathord{-}|\longrightarrow}\deleted{\unsure{X}}M\times\deleted{\unsure{X}}M$} on trouve un foncteur \deleted{$X\longrightarrow M$}
\[
 X\longrightarrow M\times M
\]
i.e. un objet \inserted{$(F_0,F_1)$} de $M(X)\times M(X)$, \deleted{d'où \unsure{une|un} \unsure{dé.}} je dis que \deleted{\unsure{ce sont deux|ce sont les}} $F_0$ et $F_1$ sont resp. la source et le but du chemin $c$. De façon un peu plus précise, on a commutativité (stricte) dans le diagramme de foncteurs
\diagram{\begin{tikzcd}[ampersand replacement=\&,column sep=large,row sep=large]
 \underline{\mathrm{Ch}}_{\unsure{\Sigma|Z}_X}(M(X)) \arrow[r,"\sim"] \arrow[d,"{\sigma_{M(X)},\beta_{M(X)}}"] \& \underline{\mathrm{Hom}}^!(X,\underline{\mathrm{Ch}}_\Sigma(M)) \arrow[d,"{\text{déduit de }\sigma_M,\beta_M}"] \\
 M(X)\times M(X) \arrow[r,"\sim"] \& \underline{\mathrm{Hom}}^{\deleted{\unsure{!|?}}}(X,M\times M)
\end{tikzcd}}
où les foncteurs $\sigma$ sont les foncteurs source, les foncteurs $\beta$ \deleted{\unsure{sont}} les foncteurs but.

Si on se donne arbitrairement les objets $F_0,F_1$ dans $M(X)$, donc un foncteur
\[
 X\longrightarrow M\times M
\]
$\underline{\mathrm{Ch}}_{\Sigma_X}(\deleted{\unsure{d}}F_0,F_1)$ s'identifie donc à la \deleted{\unsure{sous}} \inserted{catégorie} des foncteurs \inserted{``admissibles'' \inserted{i.e. se factorisant par un $\underline{\mathrm{Ch}}_{\Sigma,n}$}} $X\xrightarrow{\unsure{c|C}}\underline{\mathrm{Ch}}_\Sigma(M)$ qui \uline{\unsure{relèvent|relient}}
% archive sheet number: 32
